\documentclass{article}

% packages
\usepackage{amsmath, amsthm, thmtools, amsfonts, amssymb, luacode, catchfile, hyperref, ifthen, graphicx, bbm}
\ifcsname c@kobocompile\endcsname
	\usepackage[a5paper, total={1072pt, 1448pt}, margin=10pt, includeheadfoot]{geometry} % set page margins
\else
	\usepackage[a4paper, margin=50pt, includeheadfoot]{geometry}
\fi
\usepackage[shortlabels]{enumitem}
\usepackage[skip=3pt, indent=0pt]{parskip}

% language
\usepackage[bidi=basic, layout=tabular, provide=*]{babel}
\ifcsname c@english\endcsname
	\babelprovide[main, import]{english}
\else
	\babelprovide[main, import, alph=letters.gershayim]{hebrew}
	\setlist[enumerate,2]{label={\alph*.}}
	\babelprovide{rl}
\fi
\babelfont{rm}[Renderer=Harfbuzz]{Libertinus Serif}
\babelfont{sf}{Libertinus Sans}
\babelfont{tt}{Libertinus Mono}

% style
\ifcsname c@stoppagebreaks\endcsname
\else
	\AddToHook{cmd/section/before}{\clearpage}	% Add line break before section
\fi
\linespread{1.3}
\setcounter{secnumdepth}{0}		% Remove default number tags from sections, this won't do well with theorems
\AtBeginDocument{\setlength{\belowdisplayskip}{3pt}}
\AtBeginDocument{\setlength{\abovedisplayskip}{3pt}}
\graphicspath{ {../images/} }

% operators
\DeclareMathOperator\cis{cis}
\DeclareMathOperator\Sp{Span}
\DeclareMathOperator\tr{tr}
\DeclareMathOperator\im{Im}
\DeclareMathOperator\re{Re}
\DeclareMathOperator\diag{diag}
\DeclareMathOperator*\lowlim{\underline{lim}}
\DeclareMathOperator*\uplim{\overline{lim}}
\DeclareMathOperator\rng{rng}
\DeclareMathOperator\Sym{Sym}
\DeclareMathOperator\Arg{Arg}
\DeclareMathOperator\Log{Log}
\DeclareMathOperator\dom{dom}
\DeclareMathOperator\supp{Supp}
\DeclareMathOperator\var{Var}
\DeclareMathOperator\cov{Cov}
\DeclareMathOperator\End{End}
\DeclareMathOperator\id{id}
\DeclareMathOperator\cl{cl}
\DeclareMathOperator\gal{Gal}
\DeclareMathOperator\dist{dist}
\DeclareMathOperator\form{form}
\DeclareMathOperator\term{term}
\DeclareMathOperator\sent{sent}
\DeclareMathOperator\Th{Th}
\DeclareMathOperator\Hom{Hom}
\DeclareMathOperator\Aut{Aut}
\DeclareMathOperator\Char{char}
\DeclareMathOperator\vol{vol}
\DeclareMathOperator\GL{GL}
\DeclareMathOperator\SL{SL}
\DeclareMathOperator\lcm{lcm}

% commands
\newcommand{\cupdot}{\mathbin{\mathaccent\cdot\cup}}
\newcommand{\indicator}{\mathbbm{1}}
\newcommand{\AAA}[0]{\mathbb{A}}
\newcommand{\Aa}[0]{\mathcal{A}}
\newcommand{\BB}[0]{\mathbb{B}}
\newcommand{\Bb}[0]{\mathcal{B}}
\newcommand{\CC}[0]{\mathbb{C}}
\newcommand{\Cc}[0]{\mathcal{C}}
\newcommand{\DD}[0]{\mathbb{D}}
\newcommand{\Dd}[0]{\mathcal{D}}
\newcommand{\EE}[0]{\mathbb{E}}
\newcommand{\FF}[0]{\mathbb{F}}
\newcommand{\Ff}[0]{\mathcal{F}}
\newcommand{\Ii}[0]{\mathcal{I}}
\newcommand{\Gg}[0]{\mathcal{G}}
\newcommand{\HH}[0]{\mathbb{H}}
\newcommand{\Hh}[0]{\mathcal{H}}
\newcommand{\Kk}[0]{\mathcal{K}}
\newcommand{\Ll}[0]{\mathcal{L}}
\newcommand{\Mm}[0]{\mathcal{M}}
\newcommand{\NN}[0]{\mathbb{N}}
\newcommand{\Nn}[0]{\mathcal{N}}
\newcommand{\PP}[0]{\mathbb{P}}
\newcommand{\Pp}[0]{\mathcal{P}}
\newcommand{\QQ}[0]{\mathbb{Q}}
\newcommand{\RR}[0]{\mathbb{R}}
\newcommand{\Rr}[0]{\mathcal{R}}
\newcommand{\Ss}[0]{\mathcal{S}}
\newcommand{\TT}[0]{\mathbb{T}}
\newcommand{\Tt}[0]{\mathcal{T}}
\newcommand{\UU}[0]{\mathbb{U}}
\newcommand{\Uu}[0]{\mathcal{U}}
\newcommand{\Vv}[0]{\mathcal{V}}
\newcommand{\Ww}[0]{\mathcal{W}}
\newcommand{\Xx}[0]{\mathcal{X}}
\newcommand{\Yy}[0]{\mathcal{Y}}
\newcommand{\ZZ}[0]{\mathbb{Z}}
\newcommand{\Zz}[0]{\mathcal{Z}}
\newcommand{\acts}[0]{\circlearrowright}
\newcommand{\explain}[2] {
	\begin{flalign*}
		 && \text{#2} && \text{#1}
	\end{flalign*}
}
\newcommand{\maketitleprint}[1][green]{
	\begin{center}
		\includegraphics[width=6cm]{dragon_#1}
	\end{center}
}

% theorem commands
\newtheoremstyle{c_remark}
	{}	% Space above
	{}	% Space below
	{}% Body font
	{}	% Indent amount
	{\bfseries}	% Theorem head font
	{}	% Punctuation after theorem head
	{.5em}	% Space after theorem head
	{\thmname{#1}\thmnumber{ #2}\thmnote{ \normalfont{\text{(#3)}}}}	% head content
\newtheoremstyle{c_definition}
	{3pt}	% Space above
	{3pt}	% Space below
	{}% Body font
	{}	% Indent amount
	{\bfseries}	% Theorem head font
	{}	% Punctuation after theorem head
	{.5em}	% Space after theorem head
	{\thmname{#1}\thmnumber{ #2}\thmnote{ \normalfont{\text{(#3)}}}}	% head content
\newtheoremstyle{c_plain}
	{3pt}	% Space above
	{3pt}	% Space below
	{\itshape}% Body font
	{}	% Indent amount
	{\bfseries}	% Theorem head font
	{}	% Punctuation after theorem head
	{.5em}	% Space after theorem head
	{\thmname{#1}\thmnumber{ #2}\thmnote{ \text{(#3)}}}	% head content

\ifcsname c@english\endcsname
	\theoremstyle{plain}
	\newtheorem{theorem}{Theorem}[section]
	\newtheorem*{theorem*}{Theorem}
	\newtheorem{lemma}[theorem]{Lemma}
	\newtheorem{proposition}[theorem]{Proposition}
	\newtheorem*{proposition*}{Proposition}

	\theoremstyle{definition}
	\newtheorem{definition}[theorem]{Definition}
	\newtheorem*{definition*}{Definition}
	\newtheorem{example}{Example}[section]
	\newtheorem{exercise}{Exercise}[section]

	\theoremstyle{remark}
	\newtheorem*{remark}{Remark}
	\newtheorem*{solution}{Solution}
	\newtheorem{corollary}[theorem]{Corollary}
	\newtheorem{notation}[theorem]{Notation}
\else
	\theoremstyle{c_plain}
	\newtheorem{theorem}{משפט}[section]
	\newtheorem*{theorem*}{משפט}
	\newtheorem{lemma}[theorem]{למה}
	\newtheorem{proposition}[theorem]{טענה}
	\newtheorem*{proposition*}{טענה}

	\theoremstyle{c_definition}
	\newtheorem{definition}[theorem]{הגדרה}
	\newtheorem*{definition*}{הגדרה}
	\newtheorem{example}{דוגמה}[section]
	\newtheorem{exercise}{תרגיל}[section]

	\theoremstyle{c_remark}
	\newtheorem*{remark}{הערה}
	\newtheorem*{solution}{פתרון}
	\newtheorem{corollary}[theorem]{מסקנה}
	\newtheorem{notation}[theorem]{סימון}
\fi

% Questions related commands
\newcounter{question}
\setcounter{question}{1}
\newcounter{sub_question}
\setcounter{sub_question}{1}

\ifcsname c@english\endcsname
	\newcommand{\question}[1][0]{
		\ifthenelse{#1 = 0}{}{\setcounter{question}{#1}}
		\section{Question \arabic{question}}
		\addtocounter{question}{1}
		\setcounter{sub_question}{1}
	}

	\newcommand{\subquestion}[1][0]{
		\ifthenelse{#1 = 0}{}{\setcounter{sub_question}{#1}}
		\subsection{Part \alph{sub_question}}
		\addtocounter{sub_question}{1}
	}
\else
	\newcommand{\question}[1][0]{
		\ifthenelse{#1 = 0}{}{\setcounter{question}{#1}}
		\section{שאלה \arabic{question}}
		\addtocounter{question}{1}
		\setcounter{sub_question}{1}
	}

	\newcommand{\subquestion}[1][0]{
		\ifthenelse{#1 = 0}{}{\setcounter{sub_question}{#1}}
		\subsection{סעיף \localecounter{letters.gershayim}{sub_question}}
		\addtocounter{sub_question}{1}
	}
\fi

% import lua and start of document
\directlua{common = require ('../common')}

\GetEnv{AUTHOR}

% headers
\author{\AUTHOR}
\date\today

\title{פתרון מטלה 02 --- יסודות בתורת ההצגות}

\DeclareMathOperator\reg{reg}
\DeclareMathOperator\Ann{Ann}

% chktex-file 44

\begin{document}
\maketitle
\maketitleprint[red]

\question{}
יהי $n \ge 2$, נראה שעבור חבורת הסימטריה $S_n$ יש בדיוק שתי הצגות חד־מימדיות לא שקולות מעל $\CC$.
\begin{proof}
	תהיינה הצגות $\rho_1, \rho_2 : S_n \to \CC \setminus \{ 0 \}$, אז קיים איזומורפיזם $T : \CC \to \CC$ ביניהן אם $T \circ \rho_2 = \rho_1 \circ T$.
	נבחין גם כי אם $\sigma \in S_n$ אז $\rho_1(\sigma) x = \alpha_{\sigma} x$ לאיזשהו ערך $\alpha_{\sigma} \in \CC \setminus \{ 0 \}$.
	נסמן באופן דומה $\rho_2(\sigma) x = \beta_{\sigma} x$ ונקבל,
	\[
		T \alpha_{\sigma}
		= \beta_{\sigma} T
	\]
	אך כמובן $T$ גם היא מכפלה בערך סקלרי מרוכב ונזהה אותה כאחת כזו, ובפרט היא חילופית, כלומר $\alpha_{\sigma} = \beta_{\sigma}$ לכל $\sigma \in S_n$.
	נובע אם כך ש־$\rho_1 \cong \rho_2 \iff \rho_1 = \rho_2$ במקרה זה.

	נסמן $\sigma = (1\ 2\ \ldots\ n), \tau = (1\ 2)$, אז $S_n = \langle \sigma, \tau \rangle$.
	אז $\tau^2 = \id$, ולכן נסיק ש־$\rho(\tau) = \pm 1$ בלבד.
	נזכור כי אפשר לבנות את $\sigma$ כהרכבה של פרמוטציות מאורך $2$, לכן $\rho(\sigma) \in \{ \pm 1 \}$ גם כן, ולבסוף נבחין כי כל החלפה נקבעת ביחידות על־ידי $\tau$ כהצמדה, ולכן $\rho(\sigma) = 1$ בלבד.
	נובע אם כך שיש בדיוק שתי הצגות אפשריות במקרה זה, אשר נקבעות על־ידי הסימן של $\sigma$.
\end{proof}

\question{}
\subquestion{}
נבנה טבלת קרקטר של $D_4$.
\begin{solution}
	נסמן $\sigma$ סיבוב ו־$\tau$ שיקוף, לכן $D_4 = \langle \sigma, \tau \rangle$.
	נתחיל בחישוב ישיר של מחלקות הצמידות של החבורה,
	\[
		\{ e \},
		\quad
		\{ \sigma, \sigma^3 \},
		\quad
		\{ \sigma^2 \},
		\quad
		\{ \tau, \tau \sigma^2 \},
		\quad
		\{ \tau \sigma, \tau \sigma^3 \}
	\]
	ולכן נובע שיש בדיוק חמש הצגות אי־פריקות של $D_4$ מעל $\CC$.
	ידועה לנו ההצגה $\rho_2 : D_4 \to \GL_2(\CC)$ הנתונה על־ידי,
	\[
		\sigma
		\mapsto \begin{pmatrix}
			0 & -1 \\
			1 & 0
		\end{pmatrix},
		\quad
		\tau
		\mapsto \begin{pmatrix}
			-1 & 0 \\
			0 & 1
		\end{pmatrix}
	\]
	ראינו כי $\rho_2$ אי־פריקה אם הקרקטר המושרה $\chi_2$ מקיים,
	\[
		\sum_{g \in D_4} \chi_2(g) \overline{\chi_2(g)}
		= |G|
	\]
	בדיקה ישירה איבר איבר מראה שאכן נקבל סכום $0$, ולכן $\rho_2$ הצגה אי־פריקה.
	ראינו כי גם מתקיים,
	\[
		\sum_{\rho \in R} {(\dim \rho)}^2
		= |G|
		= 8
	\]
	עבור $R$ קבוצת ההצגות עד כדי שקילות, ומצאנו כי $\dim \rho_2 = 2$, לכן בהכרח שאר ההצגות האי־פריקות הן ממימד $1$ בלבד.
	נסמן $\rho_1 : D_4 \to \CC^\times$ ההצגה הטריוויאלית,
	ובאופן דומה לשאלה הקודמת נוכל להסיק שהצגות ממימד $1$ שקולות אם הן שוות, לכן נוכל להסיק שההצגות האחרות מתקבלות על־ידי.
	$\tau^2 = e$ ולכן איבר זה נשלח לאיבר מסדר $2$, וכן $\tau \sigma \tau = \sigma^{-1} \implies \sigma^2 = 1$ גם כן במקרה של ההצגה (במימד $1$ היא חילופית) ולכן,
	\[
		\rho_3
		= \begin{cases}
			\sigma \mapsto -1 \\
			\tau \mapsto 1
		\end{cases},
		\quad
		\rho_4
		= \begin{cases}
			\sigma \mapsto -1 \\
			\tau \mapsto -1
		\end{cases},
		\quad
		\rho_5
		= \begin{cases}
			\sigma \mapsto 1 \\
			\tau \mapsto -1
		\end{cases}
	\]
	ונשאר לחשב את הטבלה על־ידי הצבה:
	\begin{otherlanguage}{english}
		\begin{center}
			$\begin{array}{c | c c c c c}
				X & \{ e \} & \{ \sigma, \sigma^3 \} & \{ \sigma^2 \} & \{ \tau, \tau \sigma^2 \} & \{ \tau \sigma, \tau \sigma^3 \} \\
				\rho_1 & 1 & 1 & 1 & 1 & 1 \\
				\rho_2 & 2 & 0 & -2 & 0 & 0 \\
				\rho_3 & 1 & -1 & 1 & 1 & -1 \\
				\rho_4 & 1 & -1 & 1 & -1 & 1 \\
				\rho_5 & 1 & 1 & 1 & -1 & -1
			\end{array}$
		\end{center}
	\end{otherlanguage}
\end{solution}

\subquestion{}
נבנה את טבלת הקרקטרים של $A_4$.
\begin{solution}
	הפעם נבחין כי מחלקות הצמידות הן,
	\[
		\{ (1)\},
		\quad
		\{ (1\ 2\ 3), (2\ 4\ 3), (1\ 4\ 2), (1\ 3\ 4)\},
		\quad
		\{ (1\ 3\ 2), (1\ 2\ 4), (1\ 4\ 3), (2\ 3\ 4)\},
		\quad
		\{ (1\ 2)(3\ 4), (1\ 3)(2\ 4), (1\ 4)(2\ 3)\} 
	\]
	ונסמן נציגים $e, \sigma_1, \sigma_2, \sigma_3$ בהתאמה. \\
	לכן יש בדיוק ארבע הצגות אי־פריקות, ונזכור כי גם,
	\[
		\sum_{i = 1}^4 {(\dim \rho_i)}^2
		= |G|
		= 12
		= 1^2 + 3^2 + 1^2 + 1^2
	\]
	כאשר זהו הפירוק היחיד לריבועי טבעיים מחישוב ישיר.
	נסיק שיש הצגה טריוויאלית $\rho_1(\sigma) = 1$, הצגה $\rho_2$ ממימד $3$, ועוד שתי הצגות חד־מימדיות.
	נגדיר את $\rho_2 : A_4 \to \GL_3(\CC)$ להיות ההצגה הסטנדרטית של $S_4$ מצומצמת ל־$A_4$, ברור כי היא משמרת את אי־הפריקות שלה, וידוע כי $\dim \rho_2 = 3$ כפי שרצינו.
	נבחין כי אם $\rho : A_4 \to \CC$ אז מ־$\sigma_1^3 = 1$ נסיק ${(\rho(\sigma_1))}^3 = 1$, וכן $\sigma_1 \sigma_2 = e$ ולכן $\sigma_2$ נקבע על־ידי $\sigma_1$,
	ובסוף גם $\sigma_1 \sigma_3$ צמוד ל־$\sigma_1$ ולכן גם נקבע על־ידי $\rho(\sigma_1)$, נסיק שאם $\omega$ שורש יחידה פרימיטיבי מסדר $3$ אז $\rho_3(\sigma_1) = \omega$ ו־$\rho_4(\sigma_1) = \omega^2$. \\
	נעבור לחישוב הטבלה תוך שימוש בעובדה שסכום כל שורה הוא $4$,
	\begin{otherlanguage}{english}
		\begin{center}
			$\begin{array}{c | c c c c}
				& e & \sigma_1 & \sigma_2 & \sigma_4 \\
				\rho_1 & 1 & 1 & 1 & 1 \\
				\rho_2 & 3 & 0 & 0 & -1 \\
				\rho_3 & 1 & \omega & \omega^2 & 1 \\
				\rho_4 & 1 & \omega^2 & \omega & 1
			\end{array}$
		\end{center}
	\end{otherlanguage}
\end{solution}

\question{}
תהי $G = \GL_2(\CC)$.

\subquestion{}
נמצא אינסוף הצגות חד־מימדיות לא שקולות של $G$ על המרוכבים.
\begin{solution}
	נזכור כי $\det : \GL_2(\CC) \to \CC^\times$ הוא הומומורפיזם ולכן ההצגה $\rho : \GL_2(\CC) \to \CC^\times$ המתקבלת על־ידי $\rho = \det$ היא הצגה.
	נזכור כי $\CC^2 \to \CC$ המוגדר על־ידי $c \mapsto c^n$ הוא איזומורפיזם חבורתי של $\CC$ ולכן נוכל להגדיר את,
	\[
		\rho_n : \GL_2(\CC) \to \CC^\times,
		\qquad
		\rho_n(M) = {(\det M)}^n
	\]
	לכל $n > 1$ ונקבל הצגה ייחודית.
\end{solution}

\subquestion{}
יהי $m \in \NN$ ונסמן $V_m \le \CC[x, y]$ את תת־המרחב הווקטורי ממימד $m$ הנפרש על־ידי $V_m = \Sp\{ x^0 y^m, x^1 y^{m - 1}, \ldots, x^m y^0 \}$.
לכל $g = (\begin{smallmatrix} a & b \\ c & d \end{smallmatrix}) \in G$ ו־$f \in V_m$ נגדיר,
\[
	g . f
	= f(a x + c y, b x + d y)
	\in \CC[x, y]
\]
נראה שפעולה זו מגדירה הצגה $(m + 1)$ מימדית של $G$ על $V_m$.
\begin{proof}
	נתחיל ונראה שהפעולה שהגדרנו היא אופרטור, כלומר $g . f \in V_m$ לכל $g, f$.
	נבחין כי מהגדרת $V_m$ מתקבל בדיוק ש־${(x + y)}^m \in V_m$.
	יהי $\alpha x^i y^{m - i} \in V_m$ מונום, אז מתקיים,
	\[
		g . (\alpha x^i y^{m - i})
		= \alpha {(a x + c y)}^i {(b x + d y)}^{m - i}
	\]
	נבחין כי זהו פולינום כך שכל מונום בו הוא מהצורה $\alpha a^k c^{i - k} x^i y^{k - i} b^l d^{m - i - l} x^{l} y^{m - i - l}$ ולאחר בדיקת סכום חזקות $x, y$ נקבל שאכן $g . f \in V_m$.
	ברור כי חיבור מונומים נשמר ב־$V_m$, ולכן קיבלנו שזהו אכן אופרטור.

	נבחין בנוסף שמתקיים $I . f = f(1 \cdot x + 0 \cdot y, 0 \cdot x + 1 \cdot y) = f(x, y) = f$, כלומר $I \in G$ משמר זהות,
	וכדי להראות שזוהי הצגה נראה שגם $(g h) . f = g . (h . f)$, כלומר שזהו אכן הומומורפיזם.
	נסמן $h = (\begin{smallmatrix} \alpha & \beta \\ \gamma & \delta \end{smallmatrix})$ ולכן,
	\[
		g h
		= \begin{pmatrix}
			a \alpha + b \gamma & a \beta + b \delta \\
			c \alpha + d \gamma & c \beta + d \delta
		\end{pmatrix}
		\implies
		(g h) . f
		= f( (a \alpha + b \gamma) x + (c \alpha + d \gamma) y, (a \beta + b \delta) x + (c \beta + d \delta) y)
	\]
	מהצד השני,
	\[
		g . (h . f)
		= g . f(\alpha x + \gamma y, \beta x + \delta y)
	\]
	ושוב על־ידי פתיחת מונומים ובדיקה ישירה נקבל שוויון.
\end{proof}

\subquestion{}
ידוע כי $V_m$ היא העתקה אי־פריקה לכל $m$, נראה זאת עבור $m = 1$ ו־$m = 2$.
\begin{proof}
	נבחין כי $V_1 = \Sp\{ x, y \}$, כלומר כלל הפונקציות הלינאריות בשני משתנים, לכן ניתן לכתוב כל $f \in V_1$ כ־$f = A {(x, y)}^t$ עבור $A \in M_{2 \times 1}(\CC)$.
	יהי תת־מרחב $U = \Sp\{ f \} \le V_1$, ונבחין כי $g = (\begin{smallmatrix} 0 & 1 \\ 1 & 0 \end{smallmatrix}) \in G$ מקיימת $g . f(x, y) = f(y, x)$, אז אם $f = a x + b y$ עבור $x \ne y$ סיימנו.
	אם $f = a (x + y)$ נוכל לבחור את $h = (\begin{smallmatrix} -1 & 0 \\ 0 & 1 \end{smallmatrix})$ ונקבל $h . f = a(-x + y)$ ולכן $f \notin U$.
	בכל המקרים מצאנו שקיימת פעולה שלא משמרת אינווריאנטיות של תת־מרחבים ממש של $V_1$ ולכן ההצגה היא אי־פריקה.

	נעבור למקרה של $V_2 = \Sp\{ x^2, x y, y^2 \}$.
	יהי $U < V_2$, אם $\dim U = 1$ אז בדרך שקולה לחלק הקודם נוכל להראות ש־$U$ לא אינווריאנטי.
	גם במקרה זה $g$ מקיימת את הדרוש עבור מרחבים לא סימטריים ושוב $h$ מפרקת מרחבים סימטריים.
\end{proof}

\question{}
נסמן ב־$M_n(D)$ את חוג המטריצות $n \times n$ מעל חוג החילוק $D$.

\subquestion{}
נראה ש־$M_n(D)$ הוא חוג פשוט, כלומר שאין לו אידאלים לא טריוויאליים.
\begin{proof}
	יהי $I \triangleleft M_n(D)$ אידאל כלשהו. אם $I = (0)$ אז הוא טריוויאלי וסיימנו, אחרת תהי $M \in M_n(D) \setminus \{ 0 \}$ מטריצה כלשהי.
	על־ידי סגירות לכפל במטריצות החלפת שורות וסגירות של $I$, ויחד עם העובדה שאם $a \in D^\times$ אז הוא הפיך, כלומר קיימת מטריצת כפל שורה בסקלר ששולחת ליחידה נוכל לבצע הליך שקול ללכסון ולקבל $I_n \in M$,
	כלומר $I = (I_n) = M_n(D)$.
\end{proof}

\subquestion{}
נמצא את כל תת־המודולים של $M_n(D)$ כ־$M_n(D)$־מודולים, באופן שקול אלו הם הם האידאלים השמאליים.
\begin{solution}
	נזכור כי אידאל שמאלי הוא תת־חבורה חיבורית סגור לכפל באיברי החוג משמאל, אם $I \triangleleft M_n(D)$ ו־$i \in I$ אז לכל $r \in M_n(D)$ מתקיים $r \cdot i \in I$.
	כהמשך של הסעיף הקודם ברור כי אם $M \in I$ אז גם כל מטריצה שקולת שורה גם היא נמצאת ב־$I$, וכך גם ישנה סגירות לחיבור שורות ואיפוסן, אך לא לשינוי ערכי שורה.
\end{solution}

\question{}
יהי $R$ חוג ו־$M$ ו־$N$ $R$־מודולים (שמאליים).

\subquestion{}
נגדיר את $\Hom_R(M, N)$ כאוסף כל הפונקציות $\varphi : M \to N$ כך ש־$\varphi$ הומומורפיזם של חבורות אבליות וכן $\varphi(r . m) = r . \varphi(m)$ לכל $r \in R, m \in M$. \\
נראה ש־$\Hom_R(M, N)$ היא חבורה אבלית.
\begin{proof}
	נראה שזו חבורה.
	נסמן $\phi \in \Hom_R(M, N)$ העתקת האפס, כלומר $\phi(x) = 0_N$ לכל $x$.
	נראה שזוהי חבורה אבלית עם $\phi$ כאפס.
	תהי $\varphi \in \Hom_R(M, N)$ העתקה, אז מתקיים,
	\[
		(\varphi + \phi)(x)
		= \varphi(x) + \phi(x)
		= \varphi(x)
	\]
	כשכמובן הטענה נכונה גם לחיבור בכיוון ההפוך.
	נזכור ש־$N$ מודול ולכן חבורה אבלית, אז ל־$\varphi(x)$ קיים הופכי חיבורי ונסמן $-x$.
	נגדיר את $\varphi^{-1}(x) = -x$ לכל $x$, אז $\varphi^{-1} \in \Hom_R(M, N)$ מבדיקה ישירה וכן,
	\[
		(\varphi + \varphi^{-1})(x)
		= \varphi(x) - \varphi(x)
		= \phi(x)
	\]
	ונסיק שזהו הופכי ביחס לפעולה, ובהתאם $\Hom_R(M, N)$ חבורה, האבליות נובעת מהמהלכים שעשינו עד כה.
\end{proof}

\subquestion{}
נוכיח שאם $R$ קומוטטיבי אז $(r . \varphi)(m) = \varphi(r . m)$ מגדיר $R$־מודול על $\Hom_R(M, N)$.
\begin{proof}
	בסעיף הקודם הוכחנו כי זו חבורה אבלית, לכן עתה מספיק לבדוק שההעתקה $\rho : R \to \End(\Hom_R(M, N))$ המוגדרת על־ידי $(\rho(r)(\varphi))(m) = \varphi(r . m)$ היא הומומורפיזם חוגים,
	כלומר עלינו להראות ש־$r . (\varphi + \psi) = (r . \varphi) + (r . \psi)$ וכן $r . (\varphi \circ \psi) = (r \circ \varphi) \circ (r \circ \psi)$, וכן שימור יחידה כפלית.
	שימור יחידה כפלית נובע ישירות מהגדרה, ושימור חיבורי נובע מסעיף א' והגדרה, נראה שימור פעולת כפל.
	מתקיים $(r . (\varphi \circ \psi))(x) = \varphi(r . \psi(x)) = \varphi(\psi(r . x))$ כפי שרצינו בדיוק.
\end{proof}

\question{}
יהי $M$ $R$־מודול, המשמיד של $M$ מוגדר על־ידי $\Ann_R(M) = \{ r \in R \mid \forall m \in M,\ r . m = 0 \}$. \\
נראה ש־$\Ann_R(M)$ הוא אידאל של $R$.
\begin{proof}
	יהיו $r, s \in \Ann_R(M)$,
	\[
		(r + s) . m
		= r . m + s . m
		= 0 + 0
		= 0
	\]
	מתכונות הומומורפיזם, קיבלנו שהמשמיד סגור לחיבור.
	בנוסף ברור כי $0 . m = 0$ לכל $m$ ולכן $0 \in \Ann_R(M)$.
	באופן דומה אם $r \in \Ann_R(M)$ אז $(-r) . m = r . (-m) = 0$ לכל $m \in M$, ולכן $\Ann_R(M)$ חבורה חיבורית.

	נראה סגירות לכפל.
	יהי $r \in \Ann_R(M)$ ויהי $s \in R$, וכן $m \in M$.
	מתקיים,
	\[
		(s \cdot r) . m
		= s . (r . m)
		= s . 0
		= 0
	\]
	כאשר $r . m = 0$ משייכות למשמיד, לכן מצאנו סגירות משמאל.
	נראה סגירות מימין,
	\[
		(r \cdot s) . m
		= r . (s . m)
		= r . n
		= 0
	\]
	כאשר $n \in M$ ושוב מהעובדה ש־$r . m = 0$ לכל $m$.
\end{proof}

\question{}
יהי $\FF$ שדה ונבחן את $\FF[x]$.

\subquestion{}
נגדיר ונראה התאמה בין $\FF[x]$־מודולים לבין זוגות $(V, T)$ של מרחבים וקטוריים $V$ מעל $\FF$ והעתקות לינאריות $T : V \to V$.
\begin{proof}
	יהי $M$ $\FF[x]$־מודול.
	נגדיר את מרחב הפולינומים $V = M$ ונגדיר,
	\[
		\forall \alpha \in \FF, m \in M,\ 
		\alpha . m
		= (\alpha x^0) . m
	\]
	ונקבל מבנה של מרחב וקטורי.
	נגדיר את ההעתקה הלינארית $T : M \to M$ על־ידי $T(m) = x . m$, אז ברור כי היא מתחלפת עם כפל בסקלר כפולינום מעל $\FF$, ולכן זוהי העתקה לינארית.
	ברור כי $(V, T)$ היא התאמה חד־חד ערכית עד כדי איזומורפיזם מטעמי שוויון.

	יהי זוג $(V, T)$ ונגדיר את $V$ כ־$\FF[x]$־מודול על־ידי הפעולה $x . v = T(v)$. ברור כי $V$ הוא חבורה אבלית ביחס לחיבור, ולכן נשאר להראות שהפעולה הנוצרת על־ידי $x . v$ מקיימת את התכונות השקולות למודול.
	נבחין כי $x^n . v = T^n(v)$ וכן $(\alpha x^n) . v = \alpha T^n(v)$, עלינו להראות כי זהו הומומורפיזם של חוגים.
	עבור חיבור נזכור כי העתקה לינארית משמרת חיבור ולכן נגדיר,
	\[
		(\alpha x^n + \beta x^m) . v
		= \alpha T^n(v) + \beta T^m(v)
		= (\alpha x^n) . v + (\beta x^m) . v
	\]
	ונקבל שהחיבור משתמר.
	נותר לבדוק שכפל מונומים נשמר,
	\[
		(x^n x^m) . v
		= T^{n + m}(v)
		= T^n(T^m(v))
		= x^n . x^m . v
	\]
	ולכן זהו אכן הומומורפיזם ומצאנו $\FF[x]$־מודול. \\
	נבחין כי המודול המתאים הוא חד־חד ערכי לבחירת מרחב וקטורי עד כדי איזומורפיזם ישירות משוויון.
\end{proof}

\subquestion{}
בהינתן $\FF[x]$־מודול $V$, נראה שקיימת התאמה הפיכה בין תת־מודולים של $V$ ותת־מרחבים וקטוריים $W \le V$ המקיימים $T(W) \le W$.
\begin{proof}
	מהסעיף הקודם קיימת העתקה לינארית יחידה $T : V \to V$ כך ש־$(V, T)$ מתאימה ל־$V$ כ־$\FF[x]$־מודול.
	יהי תת־מודול $W \le V$ ונראה שהוא תת־מרחב וקטורי וש־$T$ אינווריאנטית אליו.
	מהגדרת תת־מודול הוא סגור חיבורית ולכפל ב־$\FF[x]$, כלומר $W$ סגורה לחיבור ולכפל הן בסקלר והן ב־$T$ כפי שהגדרנו אותה בסעיף הקודם, בפרט זהו אכן תת־מרחב וקטורי ו־$T$ אינווריאנטית אליו.

	מהצד השני יהי $W \le V$ תת־מרחב וקטורי כך ש־$T(W) \le W$.
	נבחין כי $W$ סגור לחיבור, כפל בסקלר וכפל ב־$x$, ולכן סגור לכפל בפולינום ובהתאם תת־מודול.
\end{proof}

\question{}
יהי $R$ חוג.

\subquestion{}
נראה ש־$R$ הוא גם $R$־מודול המוגדר על־ידי $r . s = r s$.
\begin{proof}
	עלינו להראות ש־$R$ הוא חבורה חיבורית, אך הוא חוג ובפרט חבורה חיבורית, וכן שההעתקה $\rho : R \to \End(R)$ המוגדרת על־ידי $\rho(r)(s) = r s$ היא הומומורפיזם חוגים.
	מהגדרה $1 . s = s$ לכל $s \in R$ ולכן יש שימור יחידה.
	בנוסף גם $(r_1 r_2) . s = (r_1 r_2) s = r_1 (r_2 s) = r_1 . (r_2 . s)$, כלומר כפל נשמר, ובמהלך זהה לחלוטין גם חיבור נשמר וקיבלנו את המבוקש.
\end{proof}

\subquestion{}
נסביר למה אנדומורפיזם חוגים $R \to R$ הוא שונה מאנדומורפיזם $R \to R$ ב־$R$־מודול.
נמצא דוגמה לאנדומורפיזם בכל קטגוריה שלא נמצא בקטגוריה השנייה.
\begin{solution}
	אנדומורפיזם חוגים משמר את פעולת הכפל של החוגים, לעומת זאת אנדומורפיזם מודולים משמר רק חיבור (כחבורה אבלית), ומשמר כפל במובן הסקלרי.

	כדי להבהיר את הדקות נבחן את $R = \ZZ$ ואת שתי ההעתקות,
	\[
		f, g : R \to R,
		\quad
		f(n) = n,
		\quad
		g(n) = 2 n
	\]
	ראינו כי $f$ היא ההומומורפיזם היחיד $\ZZ \to \ZZ$ (ולא נתייחס להומומורפיזם האפס), לעומת זאת $g$ היא ההעתקה המושרית על־ידי $2 . f \in \End_\ZZ(\ZZ)$ ובהתאם היא מציגה את האופן בו סגירות לכפל בסקלר שונה מסגירות לכפל.

	בכיוון ההפוך נבחן את $R = \FF[x]$ לאיזשהו שדה כך ש־$\Char \FF \ne 2$, ואת אנדומורפיזם ההצבה $f : R \to R$ המוגדר על־ידי $f(x) = -x$.
	ראינו כי זה אכן אנדומורפיזם של חוגי פולינומים, ונרצה להראות שזהו לא אנדומורפיזם מודולים.
	שימור של פעולת חיבור כמובן נשמר, לכן נבדוק את פעולת הכפל בסקלר.
	נניח בשלילה שזהו אנדומורפיזם מודולים ונקבל $x . (-x) = x . f(x) = f(x^2) = x^2$ בסתירה.
\end{solution}

\end{document}
